% Created 2026-08-25 mar 19:28
% Intended LaTeX compiler: pdflatex
\documentclass[11pt]{article}
\usepackage[utf8]{inputenc}
\usepackage[T1]{fontenc}
\usepackage{graphicx}
\usepackage{longtable}
\usepackage{wrapfig}
\usepackage{rotating}
\usepackage[normalem]{ulem}
\usepackage{amsmath}
\usepackage{amssymb}
\usepackage{capt-of}
\usepackage{hyperref}
\author{Marisol}
\date{\today}
\title{}
\hypersetup{
 pdfauthor={Marisol},
 pdftitle={},
 pdfkeywords={},
 pdfsubject={},
 pdfcreator={Emacs 29.3 (Org mode 9.6.15)}, 
 pdflang={English}}
\begin{document}

\tableofcontents

\section{Diseño de Algoritmos - TP0}
\label{sec:org0a4bca7}


Este trabajo práctico tiene como objetivo repasar conceptos fundamentales sobre eficiencia,
estructuras de datos y complejidad, así como diagnosticar el nivel de conocimientos previos
de los estudiantes.

\subsection{Parte 1: Análisis de complejidad}
\label{sec:org74505f2}

Analizar la complejidad temporal de los siguientes fragmentos de código. Justificar
utilizando reglas de suma y producto.

Ejercicio 1:

1.for (i=1; i<n; i++) \{ // O(n)

2.for (j=1; j<=n; j++) \{//O(n)

3.if (a[i] < a[j]) //O(max(if,else) + 1 )

4.resp = true;  //O(1)

5.else

6.resp = false; //O(1)

7.k = i + j; //O(1)

\}
\}

El orden total es O(n²)


Ejercicio 2:

i = 1; //O(1)

while (i <= n) \{ //O(log n)

if (a[i] >= a[n]) a[n] = a[i]; //O(1)

i = i * 2;

\}


El orden total es O(log n)




\section{Parte 2: Estructuras de datos y operaciones}
\label{sec:org407af14}

\subsection{Ejercicio 3:}
\label{sec:orgf03ebbe}

Determinar la complejidad (Big-O) de las siguientes operaciones y justificar la respuesta:

a. Insertar un elemento en la posición k de una Lista. //O(n)

b.Poner un elemento en una Cola. //O(1)

c.Buscar un determinado elemento en un Árbol Binario. 

En el peor caso considerando que el árbol no está ordenado, es de O(n)
En el mejor caso el orden sería O(log n)



d. Eliminar un elemento en un Árbol Binario de Búsqueda.

Peor caso: Teniendo en cuenta que el árbol esté desbalanceado, el peor
caso será O(n)


Mejor caso: Sí el árbol está balanceado y ordenado, el mejor caso será
O(log n)

e.Insertar un elemento en una tabla Hash. //O(1)

\subsection{Ejercicio 4:}
\label{sec:org34224d8}

Analizar las clases de API Java para manejo de estructuras de datos y sus respectivos
métodos. Tomar Queue y Hashmap como ejemplo.


\subsection{Ejercicio 5:}
\label{sec:org0bbc145}

Un juego de ajedrez se puede plantear como el siguiente problema de decisión:

Dado un posicionamiento legal de las piezas de ajedrez y la información sobre qué mueve,
determinar si ese lado puede ganar. ¿Es decidible este problema de decisión? Explicar.

Este algoritmo no es decidible, ya que no es posible determinar si por
un movimiento de pieza las piezas blancas ganarán. Las posibles
jugadas de un juego de ajedrez son 10\textsuperscript{120}.
Teoricamente sí es decidible pero aplicado a la realidad no. Si bien
tiene una cota superior finita, es imposible alcanzar ese número en un
algoritmo con la tecnología actual.



\section{Ejercicio 5:}
\label{sec:org753dd78}

Un determinado problema puede ser resuelto por un algoritmo cuyo tiempo de ejecución
está en O (n log2 n). ¿Cuál de las siguientes afirmaciones es verdadera? Explicar.

El problema es manejable.

Es manejable porque su crecimiento si bien es inclinado, es más suave
que x². 

El problema es intratable.

Imposible decirlo.



\section{Parte 3: Diagnóstico conceptual}
\label{sec:org05b545e}

a) ¿Qué entendemos por eficiencia de un algoritmo?

La eficiencia de un algoritmo la determina el tiempo, los recursos que
se utilizan para ejecutar a este y el espacio utilizado (memoria), todo esto
considerando el peor caso.

b) ¿Qué estructuras de datos conocés y usás con mayor seguridad?


LIsta, árbol, hashmap.




\section{Parte 4 (opcional): Problema abierto}
\label{sec:orge9e5733}

Dado un conjunto de números enteros positivos, diseñar un algoritmo (o describir en pasos)
que determine si es posible dividir el conjunto en dos subconjuntos cuya suma sea igual. No
es necesario programarlo, pero se valorará la claridad de la explicación y el enfoque
algorítmico.



Algoritmo:

\begin{enumerate}
\item Sumar todos los números del conjunto grande.
\item 
\end{enumerate}
\end{document}
